\chapter{Da bancada à plataforma}
\label{cap:operacao}

\begin{objetivos}
Ao final deste capítulo, você será capaz de:
\begin{itemize}[leftmargin=*,itemsep=0.2em]
  \item enumerar os requisitos não funcionais que distinguem um protótipo de um sistema
        implantado;
  \item enunciar os princípios de uma interface que comunique incerteza sem induzir
        fadiga de alarme;
  \item estruturar o laço de realimentação humana e a governança de classes virtuais;
  \item reconhecer e monitorar deriva de conceito e de dados;
  \item estimar o custo total de operação de um sistema de monitoramento com camada de
        linguagem.
\end{itemize}
\end{objetivos}

\section{O que muda quando o sistema é ligado}

Todos os números deste livro foram obtidos em avaliação controlada. A exceção é o
\cref{cap:dual}, cujo sistema opera em mais de vinte unidades flutuantes --- e é
justamente por isso que ele é o capítulo com o menor número de tabelas e o maior número
de decisões de engenharia aparentemente banais.

Este capítulo trata dessas decisões. Elas não aparecem em artigos, e determinam se o
sistema é usado ou desligado.

\begin{destaque}[A assimetria entre pesquisa e operação]
Em pesquisa, o critério é a métrica. Em operação, o critério é: \emph{o operador confia?}

Um sistema com acurácia de $0{,}95$ que emite três alarmes falsos por turno será
ignorado em uma semana. Um sistema com acurácia de $0{,}85$ que quase nunca alarma
indevidamente será consultado. A métrica não captura essa diferença; o desvio padrão
entre poços --- discutido no \cref{cap:segmentacao} --- captura parte dela.
\end{destaque}

\section{Requisitos não funcionais}

\begin{table}[htbp]
\centering
\caption{Requisitos não funcionais e a resposta arquitetural correspondente.}
\label{tab:oper-requisitos}
\small
\begin{tabular}{p{3.6cm}p{8.4cm}}
\toprule
\textbf{Requisito} & \textbf{Resposta no sistema descrito} \\
\midrule
Disponibilidade desde o primeiro segundo &
Diagrama de decisão determinístico, que não requer treinamento
(\cref{cap:dual}) \\
\addlinespace
Ausência de rotulação contínua &
Autoencoder e ciclo de mundo aberto operam sem rótulos \\
\addlinespace
Heterogeneidade entre poços &
Modelo por poço; normalização por escore-$z$; avaliação com desvio padrão entre poços
reportado \\
\addlinespace
Robustez a falha de sensor &
Camada de validação de dado; imputação; detecção de sensor congelado \\
\addlinespace
Robustez a manobra deliberada &
Reinício do ciclo de treinamento na mudança de válvula \\
\addlinespace
Latência compatível &
\SI{405}{\milli\second} do ensemble contra 30 min a 5 h de tempo de resposta do processo \\
\addlinespace
Auditabilidade &
Métricas determinísticas verificáveis; justificativas ancoradas nelas \\
\addlinespace
Extensibilidade &
Classificadores binários independentes; classes virtuais \\
\bottomrule
\end{tabular}
\end{table}

\begin{destaque}[Nenhum desses requisitos é sobre acurácia]
Vale notar o que a \cref{tab:oper-requisitos} não contém. Nenhuma linha exige acurácia
mínima.

Isso não significa que a acurácia seja irrelevante --- significa que ela é condição
necessária e largamente insuficiente. Um sistema perfeitamente acurado que exige
rotulação contínua, ou que leva uma semana para ficar disponível, ou cujo comportamento
varia imprevisivelmente entre poços, não é implantável.
\end{destaque}

\section{A interface}

\subsection{O que apresentar}

O sistema produz, para cada segmento anômalo, quatro objetos distintos:

\begin{enumerate}[leftmargin=*,itemsep=0.2em]
  \item a delimitação temporal do evento;
  \item a decisão conhecido/novidade, com o escore de Mahalanobis;
  \item a classe, quando conhecida, com o escore do classificador binário;
  \item três hipóteses ordenadas com justificativa, do agente.
\end{enumerate}

A tentação é fundir tudo em um número. O \cref{cap:closedworld} já argumentou contra, ao
discutir o painel que exibe risco analítico e risco aprendido lado a lado
(\cref{fig:closed-painel}).

\begin{destaque}[Indicadores que discordam são informativos]
Quando o método geométrico diz ``conhecido, classe 7'' e o agente diz ``incompatível com
todas as classes'', essa discordância é o achado mais valioso que o sistema pode
produzir naquele instante. Fundir os dois em um índice a destruiria.

A regra de projeto: \textbf{fundir indicadores é aceitável quando eles concordam e
inaceitável quando discordam} --- e como não se sabe de antemão qual será o caso, não se
funde.
\end{destaque}

\subsection{Graduação em vez de binário}

O \cref{cap:dual} usa três níveis de risco --- baixo, médio, alto --- em vez de um
alarme binário. A justificativa merece ser repetida: um alarme binário força uma escolha
entre perder eventos e cansar o operador, e essa escolha é feita por quem projeta o
sistema, que tem menos contexto que quem o opera.

A graduação transfere a decisão para onde há mais informação.

\subsection{Validação de agrupamentos}

\begin{figure}[htbp]
\centering
\includegraphics[width=\textwidth]{operator_mockup_interface.png}
\caption{Interface de validação de agrupamentos. O especialista inspeciona a projeção do
espaço latente, seleciona um ponto e recebe as variações percentuais por sensor daquele
segmento. Pode então reatribuir o ponto a outro agrupamento ou removê-lo da análise.}
\label{fig:oper-validacao}
\end{figure}

A \cref{fig:oper-validacao} materializa o laço humano do ciclo de mundo aberto. Três
decisões de projeto merecem comentário.

\textbf{A projeção é interativa, não estática.} O especialista precisa explorar, não
apenas contemplar. As ressalvas do \cref{cap:novidades} sobre projeções não lineares
valem aqui: a interface deve ser entendida como ferramenta de navegação, não como
representação fiel de distâncias.

\textbf{Cada ponto expõe suas estatísticas.} Selecionar um ponto revela as variações
percentuais por sensor --- as mesmas métricas do \cref{cap:agente}. O especialista julga
sobre números verificáveis, não sobre a posição visual.

\textbf{As ações são reversíveis e explícitas.} Reatribuir e remover, ambas registradas.
A remoção é particularmente importante: nem tudo que foi rejeitado merece virar classe.

\begin{destaque}[Note o valor de $15722\,\%$ na figura]
A variação percentual de T-TPT exibida é de mais de quinze mil por cento. Isso não é um
evento físico --- é o comportamento da \cref{eq:ag-pct} quando $\bar{b}$ é próximo de
zero, tipicamente porque o sensor estava congelado ou zerado no regime de base.

A interface expõe esse valor sem tratamento, e isso é uma virtude: o especialista
reconhece imediatamente um artefato de instrumentação. Uma interface que ``limpasse'' o
número o esconderia.

É também mais uma manifestação do artefato de telemetria que o
\cref{cap:resultados-agente} identificou.
\end{destaque}

\section{O laço humano}

\subsection{Três papéis distintos}

\begin{description}[leftmargin=0pt,style=nextline,itemsep=0.4em]

\item[Operador de turno.] Recebe alertas graduados, decide se investiga, registra o
desfecho. Precisa de latência baixa e de baixa taxa de falso alarme.

\item[Especialista de disciplina.] Valida agrupamentos, nomeia classes virtuais, corrige
perfis. Atua em escala de dias, não de minutos.

\item[Responsável pelo modelo.] Monitora deriva, decide retreinamentos, versiona
modelos e prompts.

\end{description}

Confundir esses papéis é um erro comum. A interface do operador não deve expor controles
de validação de agrupamento; a do especialista não precisa de latência de segundos.

\subsection{Governança de classes virtuais}

Um sistema que promove classes virtuais indefinidamente acumula categorias mal definidas.
A \cref{tab:oper-governanca} propõe uma política.

\begin{table}[htbp]
\centering
\caption{Política de governança de classes virtuais.}
\label{tab:oper-governanca}
\small
\begin{tabular}{p{3.2cm}p{8.8cm}}
\toprule
\textbf{Decisão} & \textbf{Critério} \\
\midrule
Promover a virtual &
Grupo aprovado nos três índices internos; ao menos trinta segmentos; convergência de
nomes acima de um limiar \\
\addlinespace
Encaminhar a revisão &
Toda classe virtual, dentro de um prazo definido, com o nome e a justificativa do
agente \\
\addlinespace
Fundir &
Dois grupos cujas descrições convergem para o mesmo nome e cujos centroides distam menos
que um limiar \\
\addlinespace
Descartar &
Grupo cuja dispersão de nomes indica confabulação; grupo dominado por artefato de
instrumentação \\
\addlinespace
Promover a real &
Confirmação do especialista, com definição por mecanismo e não por sintoma \\
\bottomrule
\end{tabular}
\end{table}

\begin{destaque}[A última linha é a mais importante]
``Definição por mecanismo e não por sintoma'' é a lição que as classes 4 e 5 ensinaram ao
longo de todo o livro. Uma política de governança que não a incorpore reproduzirá o
problema em cada nova classe.
\end{destaque}

\section{Deriva}

\subsection{Três formas}

\begin{description}[leftmargin=0pt,style=nextline,itemsep=0.4em]

\item[Deriva de dados.] A distribuição das entradas muda --- o poço envelhece, o BSW
sobe, a pressão de reservatório cai. O modelo continua correto sobre o passado e
progressivamente irrelevante sobre o presente.

\item[Deriva de conceito.] A relação entre entrada e rótulo muda. Um mesmo padrão de
sensores passa a significar outra coisa, por exemplo após uma intervenção que altera a
completação.

\item[Deriva induzida pelo próprio sistema.] O caso mais insidioso: o modelo se retreina
sobre dados que ele mesmo classificou como normais e, gradualmente, aprende a considerar
normal uma degradação progressiva.

\end{description}

\begin{destaque}[A salvaguarda do \cref{cap:dual}, revisitada]
A regra de excluir dados anômalos do conjunto de retreinamento existe precisamente para
impedir a terceira forma. Sem ela, uma incrustação que se desenvolve ao longo de meses
seria progressivamente absorvida como normalidade, e o alarme se extinguiria sozinho ---
sem que nenhuma métrica indicasse falha.

É o modo de falha mais perigoso de sistemas auto-adaptativos, porque é silencioso.
\end{destaque}

\subsection{O que monitorar}

\begin{itemize}[leftmargin=*,itemsep=0.3em]
  \item \textbf{Distribuição do escore de Mahalanobis} sobre dados classificados como
        normais. Um deslocamento sistemático indica deriva de dados.
  \item \textbf{Taxa de rejeição.} Um aumento sustentado pode indicar deriva ou o
        surgimento de um fenômeno novo --- e distinguir os dois é trabalho do
        especialista.
  \item \textbf{Taxa de confirmação do agente.} Queda na fração de decisões validadas
        sugere que os perfis de classe estão desatualizados.
  \item \textbf{Convergência de nomes.} Aumento de dispersão indica que o agente está
        confabulando com mais frequência.
  \item \textbf{Desfecho registrado pelo operador.} A única fonte de rótulo verdadeiro
        em operação, e a mais valiosa.
\end{itemize}

\begin{destaque}[Quatro dos cinco indicadores não exigem rótulo]
Isso é essencial. Um esquema de monitoramento que dependesse de rótulos seria inútil
exatamente onde é mais necessário --- em operação contínua, sem anotação.

O quinto indicador, o desfecho registrado, exige disciplina operacional e é o mais
difícil de obter. Vale investir nele: é o que permite avaliar o sistema de verdade.
\end{destaque}

\section{Custo}

\begin{table}[htbp]
\centering
\caption{Estrutura de custo por camada, em ordem de grandeza.}
\label{tab:oper-custo}
\small
\begin{tabular}{p{3.4cm}p{4.2cm}p{4.4cm}}
\toprule
\textbf{Camada} & \textbf{Custo de inferência} & \textbf{Observação} \\
\midrule
Diagrama de decisão & Desprezível & Regras \\
U-Net & Milissegundos por sensor & GPU opcional \\
Ensemble binário & \SI{405}{\milli\second}; \SI{21}{\mega\byte} & Por poço \\
Mahalanobis & Desprezível & Inversão de matriz $64 \times 64$ pré-computada \\
Agente de linguagem & Segundos; custo por chamada & \textbf{Somente sobre segmentos anômalos} \\
\bottomrule
\end{tabular}
\end{table}

\begin{destaque}[A decisão de custo que viabiliza a camada de linguagem]
Acionar o agente a cada janela, em duzentos e cinquenta poços, seria proibitivo em custo
e desnecessário: a esmagadora maioria das janelas é normal.

Acioná-lo apenas sobre segmentos que a camada 4 sinalizou reduz o volume em ordens de
grandeza --- de milhões de janelas por dia para dezenas de segmentos. É por isso que a
arquitetura do \cref{cap:sintese} coloca o agente na posição seis, e não na posição um.

Generalizando: em sistemas com componentes de custo muito assimétrico, a arquitetura
deve garantir que o componente caro só seja acionado quando o barato já filtrou.
\end{destaque}

\subsection{Custos que não aparecem na tabela}

\begin{itemize}[leftmargin=*,itemsep=0.3em]
  \item \textbf{Tempo do especialista} na validação de classes virtuais. É o custo
        dominante do ciclo de mundo aberto, e a razão pela qual as descrições geradas
        pelo agente têm valor econômico direto.
  \item \textbf{Versionamento de modelos e prompts.} Um prompt é código; deve ser
        versionado, revisado e testado.
  \item \textbf{Dependência de serviço externo.} Disponibilidade, latência de rede,
        mudança de modelo sem aviso.
  \item \textbf{Governança de dados.} Enviar dados operacionais a um serviço de terceiros
        exige avaliação de conformidade.
\end{itemize}

\section{Uma lista de verificação}

\begin{enumerate}[leftmargin=*,itemsep=0.3em]
  \item O sistema funciona no primeiro segundo, antes de qualquer treinamento?
  \item Ele informa explicitamente quando não sabe?
  \item Ele degrada graciosamente quando um sensor falha?
  \item Manobras deliberadas geram alarme?
  \item O desempenho foi reportado com desvio padrão entre poços?
  \item Cada afirmação do sistema pode ser rastreada a um número verificável?
  \item Indicadores discordantes são apresentados separadamente?
  \item Existe caminho para o operador registrar o desfecho?
  \item Existe política escrita de governança de classes novas?
  \item O monitoramento de deriva dispensa rótulos?
  \item O componente caro só é acionado após filtragem pelo barato?
  \item Prompts e modelos estão versionados?
\end{enumerate}

\resumocapitulo{%
A distância entre um protótipo e um sistema implantado é feita de requisitos não
funcionais que nenhuma métrica captura: disponibilidade desde o primeiro segundo,
operação sem rotulação contínua, robustez a falha de sensor e a manobra deliberada,
auditabilidade e extensibilidade. A interface deve apresentar separadamente a
delimitação, a decisão conhecido/novidade, a classe e as hipóteses do agente --- fundir
indicadores é aceitável quando concordam e inaceitável quando discordam, e como não se
sabe de antemão, não se funde. O laço humano envolve três papéis com necessidades
distintas: operador de turno, especialista de disciplina e responsável pelo modelo. A
governança de classes virtuais define quando promover, fundir, descartar e --- crucialmente
--- exige que a promoção a classe real seja feita por mecanismo, não por sintoma. Das
três formas de deriva, a mais perigosa é a induzida pelo próprio sistema, contra a qual a
exclusão de dados anômalos do retreinamento é a salvaguarda. Quatro dos cinco indicadores
de deriva propostos dispensam rótulos. Por fim, a camada de linguagem só é
economicamente viável porque é acionada apenas sobre segmentos já filtrados: em sistemas
com custos assimétricos, o componente caro deve vir depois do barato.}

\begin{exercicios}
\item Estime o volume diário de janelas processadas em duzentos e cinquenta poços com
      inferência a cada minuto, e o número esperado de segmentos anômalos com taxa de
      falso alarme de $5\,\%$. Quantas chamadas ao agente isso representa por dia?

\item Projete a interface do operador de turno com os quatro objetos produzidos pelo
      sistema. Justifique cada elemento visual em termos de carga cognitiva e de risco de
      fadiga de alarme.

\item Implemente os cinco indicadores de deriva propostos. Simule uma deriva de dados
      lenta sobre um poço do 3W e verifique quais indicadores a detectam e com quanta
      antecedência.

\item Construa um cenário concreto de deriva induzida pelo próprio sistema e demonstre,
      por simulação, que a salvaguarda de exclusão de dados anômalos a impede. O que
      acontece se a salvaguarda tiver falsos negativos?

\item A \cref{fig:oper-validacao} exibe uma variação percentual de mais de quinze mil por
      cento. Proponha três formas de apresentar esse valor ao especialista --- exibir sem
      tratamento, limitar a uma faixa, ou substituir por uma marca de artefato --- e
      discuta os riscos de cada uma.

\item Redija a política de governança de classes virtuais como documento operacional,
      com critérios numéricos concretos, prazos e responsáveis. Aplique-a
      retrospectivamente aos agrupamentos do \cref{cap:mundoaberto}.

\item Percorra a lista de verificação para um sistema de aprendizado de máquina com que
      você tenha trabalhado. Quantos itens ele satisfaz? Para cada item não satisfeito,
      estime o custo de corrigi-lo.
\end{exercicios}

\begin{leituras}
\item \citet{aranha2023system} e \citet{Aranha2024} --- descrições de sistemas de
      monitoramento em operação real.
\item \citet{LopesTese2025} --- o contexto operacional completo do sistema dual.
\item \citet{nodered} --- a ferramenta de orquestração usada na implantação.
\item \citet{gatta2022predictive} e \citet{soriano2021visual} --- aspectos operacionais e
      de visualização em monitoramento de poços.
\end{leituras}
